% R0_typecheck_repair.tex
% Build:  lualatex -interaction=nonstopmode R0_typecheck_repair.tex
\documentclass[11pt,a4paper]{article}

\usepackage{luatexja}
\usepackage{luatexja-fontspec}
\setmainjfont[BoldFont=Hiragino Sans W6]{Hiragino Sans W3}

\usepackage[margin=22mm]{geometry}
\usepackage{fontspec}
\setmonofont[Scale=0.92]{Menlo}
\usepackage{xcolor}
\usepackage{listings}
\usepackage{hyperref}
\hypersetup{colorlinks=true,linkcolor=blue!50!black,urlcolor=blue!60!black}
\usepackage{booktabs}
\usepackage{array}
\usepackage{longtable}
\usepackage{enumitem}
\usepackage{amssymb}  % \checkmark
\setlist[itemize]{leftmargin=*,topsep=2pt,itemsep=1pt}

% Code listing style
\definecolor{kw}{rgb}{0.10,0.30,0.65}
\definecolor{str}{rgb}{0.50,0.20,0.10}
\definecolor{com}{rgb}{0.30,0.55,0.30}
\definecolor{bg}{rgb}{0.97,0.97,0.95}
\lstdefinelanguage{abcl}{
  morekeywords={class,method,var,if,else,while,for,do,return,send,new,self,sender,
                int,float,string,bool,unit,any,array,function,reply,true,false,pub},
  morecomment=[l]{//},
  morestring=[b]",
}
\lstdefinelanguage{ocaml}{
  morekeywords={let,rec,in,fun,function,match,with,if,then,else,try,raise,
                module,struct,sig,end,type,of,val,open,and,as,when,for,to,do,done,
                begin,end,not,or,land,lor,lxor,lsl,lsr,asr,Some,None,Ok,Error},
  morecomment=[s]{(*}{*)},
  morestring=[b]",
}
\lstset{
  basicstyle=\ttfamily\footnotesize,
  keywordstyle=\color{kw}\bfseries,
  commentstyle=\color{com}\itshape,
  stringstyle=\color{str},
  numbers=left,numberstyle=\tiny\color{gray},
  backgroundcolor=\color{bg},
  frame=single,framerule=0.2pt,framesep=2pt,
  breaklines=true,
  showstringspaces=false,
  columns=fullflexible,
  keepspaces=true,
  xleftmargin=8pt,xrightmargin=4pt,
  belowskip=4pt,aboveskip=6pt,
}

\title{\bfseries AIPL 型チェッカー回帰修復報告 \\
       \large Xinu サンプル 3 件の修正と恒久的な型システム改善}
\author{abclcp-project R0 phase\\
        \texttt{aice-pi-evolution/experiments/2026-05-20\_xinu\_razpi\_aipl}}
\date{2026 年 5 月 20 日}

\begin{document}
\maketitle

\begin{abstract}
\noindent
本報告は、Xinu/Raspberry~Pi 進化プロジェクトの初動 phase である R0
(regression 修復) の作業内容と、その過程で明らかになった
AIPL 型チェッカー側の永続的な不具合の修復をまとめる.作業前は
\texttt{abclc/} 配下の 3 サンプル
(\texttt{Rotate4LinesXinu}, \texttt{Philosophers5Xinu},
\texttt{BoundedBufferXinu}) が全て型エラーで reject されていたが、
\texttt{src/infer.ml} の\textbf{ クラスフィールド処理に
\texttt{TypedVarDecl} を追加}し、\texttt{src/aipl2c.ml}
で\textbf{ codegen 前に \texttt{normalize\_program} を呼ぶ}、
\texttt{src/typing\_env.ml} に\textbf{ Xinu GUI 関数群と
int/float 混在の関係演算子を登録}した結果、
3 サンプルすべてが型検査・C 生成・Xinu kernel リンク・QEMU 起動の
4 段ゲートを通過するようになった.既存の非 Xinu サンプル 8 件にも
回帰は無く、合計 R0\,/\,R1\,/\,P1 で 26 個の assertion が green
になっている.
\end{abstract}

\section{背景}

\subsection{進化プロジェクトの位置づけ}

Embedded Xinu (arm-qemu / Raspberry Pi 系) 上で AIPL を進化させる
Round~1 は、5 つの方向 (R: regression 修復, P: scheduler,
F: AIPL 機能, G: GUI/I/O, N: 分散) に分かれて 21 phase + 63 assertion
を目標としている.最初のサブフェーズ\textbf{ R0} は、その出発点として
既存の Xinu サンプル 3 件が現行ツールチェーンで再び\textbf{ build + 動作}
することを保証するものである.

\subsection{2026-05-20 朝の動作確認}

aipl2c~+~xinu-raz~+~QEMU の各層は単体では動作するが、
\texttt{aipl2c --xinu} が 3 サンプル全てを\textbf{ 型検査の段階で
reject していた}ことを確認した.

\begin{lstlisting}[language=,caption={3 サンプルでの拒否メッセージ (要約)}]
$ ./_build/default/src/aipl2c.exe abclc/Rotate4LinesXinu.abcl \
      --xinu --max-msgs 0
class Spinner
  init : (int * float * float * float * int * int * int) -> 'a
  tick : () -> 'a
[Type error] line 18, col 5: type mismatch        <-- speed = sp;

$ ./_build/default/src/aipl2c.exe abclc/Philosophers5Xinu.abcl \
      --xinu --max-msgs 0
[Type error] line 68, col 16: send target is not actor: int
                                                  <-- send left_fork.request(idx)

$ ./_build/default/src/aipl2c.exe abclc/BoundedBufferXinu.abcl \
      --xinu --max-msgs 0
[Type error] line 81, col 7: type mismatch        <-- counter = 110 - slider * 10
\end{lstlisting}

\texttt{--no-typecheck} を付ければ C 生成は進むが、これは型安全網を全て
無効化してしまうため恒久的な解決にならない.

\subsection{予備調査で発覚した深層不具合}

3 ファイルの内容自体を修正しようとして、まず\textbf{ 既存の
\texttt{abclc/TypedDemo.abcl} で同じ事象}が起きていることに気付いた:

\begin{lstlisting}[language=,caption={TypedDemo.abcl ですら type-check に失敗}]
$ ./_build/default/src/aipl2c.exe abclc/TypedDemo.abcl --check
class Counter
  bump : (int) -> int
[Type error] line 37, col 13: unbound variable: count
\end{lstlisting}

\texttt{TypedDemo.abcl} は\texttt{ var count: int = 0;} という\textbf{
型注釈付きフィールド}を持ち、 \texttt{bump} メソッドがその \texttt{count}
を参照する.この pattern は AIPL の Phase~11+ で導入された
\texttt{TypedVarDecl} を用いており、ここがエラーの \emph{真の発生源}
であった.つまり、\textbf{ 型注釈付きフィールドを持つあらゆるクラスが
全 Xinu codegen 経路で壊れていた}ことになる.

\section{根本原因の分析}

\subsection{原因 1 — \texttt{infer.ml} の Class 処理から
\texttt{TypedVarDecl} が抜け落ちている}

\texttt{src/infer.ml} の \texttt{check\_decl} は、クラスフィールドを
ローカル環境に bind する際に \texttt{VarDecl} だけを認識していた:

\begin{lstlisting}[language=ocaml,caption={infer.ml の問題箇所 (修正前)}]
let check_decl (env:env) = function
  | Class c ->
    let env_cls = clone env in
    List.iter
      (fun (st:Ast.stmt) ->
        match st.sdesc with
        | VarDecl (name, init) ->
            let t   = infer_expr env_cls init in
            let sch = Types.generalize (ftv_env env_cls) t in
            set env_cls name sch
        | _ -> ()                  (* <-- TypedVarDecl がここに落ちる *)
      ) c.fields;
    ...
\end{lstlisting}

その結果\texttt{ var count: int = 0;} のフィールドは \texttt{env\_cls}
へ\textbf{ 一度も登録されない}まま methods の本文 check に入る.
method 本文中の \texttt{count = count + 1} は\textbf{ unbound variable:
count} で reject される.

\subsection{原因 2 — codegen が \texttt{normalize\_program} 抜きで動いていた}

\texttt{src/ast.ml} には全 \texttt{TypedVarDecl} を \texttt{VarDecl} に
畳んで型注釈をサイドテーブルに退避する \texttt{normalize\_program} が
存在するが、\texttt{src/aipl2c.ml} のパイプラインは型検査直後に
\textbf{normalize を経由せず}に直接 codegen を呼んでいた.
そのため C-translator 側の各 backend
(\texttt{gen\_program\_xinu}, \texttt{gen\_class}, etc.) が
field 反復で\texttt{ VarDecl} のみ拾う設計と矛盾し、
typed field は\textbf{ 列挙 (\texttt{enum F\_class\_field}) には現れる
ものの method body から \texttt{g\_field} という未宣言の global
シンボルとして参照される}という、リンク段階で必ず失敗するコードを
吐いていた.

\subsection{原因 3 — Xinu GUI 関数が gradual 化されていた}

\texttt{src/typing\_env.ml::prelude} には \texttt{sdl\_*} 系や
\texttt{ai\_call} などは登録されているが、Xinu GUI 表面の
\texttt{xinu\_gui\_set\_line} / \texttt{xinu\_gui\_slider\_value} /
\texttt{xinu\_gui\_set\_phil} など 12 個の関数は\textbf{ どれも未登録}で
\texttt{[type warning] unknown function ... treated as gradual (any -> any)}
として通っていた.この結果、

\begin{lstlisting}[language=abcl]
counter = 110 - xinu_gui_slider_value(0) * 10;
\end{lstlisting}

の \texttt{slider\_value} 戻り値は \texttt{any}、しかし \texttt{*} の
オーバーロードは \texttt{(int,int), (float,float), (int,float),
(float,int)} のみで \texttt{any} を受けるものが無く、結果
\textbf{ 型エラー (overload 失敗)} になっていた.同時に、
\texttt{if (angle >= 360)} のような\textbf{ float vs int の関係比較}も
prelude に overload が無く reject されていた.

\section{修正}

\subsection{修正 1 — \texttt{infer.ml} に \texttt{TypedVarDecl} ハンドラ}

\begin{lstlisting}[language=ocaml,caption={infer.ml::check\_decl Class — TypedVarDecl 追加}]
| TypedVarDecl (name, te, init) ->
    (* var name: T = init;  inside a class body.  Use the
       declared type T for the field's scheme so methods that
       reference the field see the annotated type.  Still
       unify against the initializer so a mismatch at the
       field declaration site surfaces immediately. *)
    let declared = ty_of_type_expr te in
    let t_init   = infer_expr env_cls init in
    if not (unify_at st.sloc declared t_init) then
      Types.type_error ~loc:st.sloc
        (Printf.sprintf "field %s: declared type does not match initializer"
           name);
    let sch = Types.generalize (ftv_env env_cls) declared in
    set env_cls name sch
\end{lstlisting}

宣言型を field の scheme として登録し、initializer は\texttt{ unify} で
追加チェックする.これで\texttt{ TypedDemo.abcl} を始めとした既存の
typed-field サンプルも自然に通る.

\subsection{修正 2 — aipl2c が \texttt{normalize\_program} を経由}

\begin{lstlisting}[language=ocaml,caption={aipl2c.ml — codegen 前に normalize}]
if !check_only then begin
  Printf.printf "[aipl2c] --check: ...\n";
  exit 0
end;
(* Strip `var name: T = e;` field/local annotations down to plain
   VarDecl (name, e) so every codegen backend sees the simpler form.
   The annotation side table is no longer consulted at codegen — only
   the type checker (which has already run) needed it. *)
let prog = Ast.normalize_program prog in
let c_code = (* gen_program_xinu / gen_program / ... *)
\end{lstlisting}

これで\texttt{ TypedVarDecl} は型検査済の AST から消え去り、Xinu
だけでなく Pony / Erlang / Go / Prolog / Python / OpenMP / LLVM
全てのバックエンドが透過的に typed-field を扱えるようになる.

\subsection{修正 3 — \texttt{typing\_env.ml} に Xinu GUI 表面 + 混在比較}

12 個の Xinu GUI 関数と、関係演算子の int/float 混在 overload を
prelude に登録した:

\begin{lstlisting}[language=ocaml,caption={typing\_env.ml — extern + mixed numeric 関係}]
add_mono e "xinu_gui_set_line"        (TFun ([TAny;TAny;TAny;TAny;TAny;TAny;TAny;TAny], TUnit));
add_mono e "xinu_gui_register_ticker" (TFun ([TAny], TUnit));
add_mono e "xinu_gui_add_button"      (TFun ([TString;TInt;TInt;TInt;TInt;TAny;TString], TUnit));
add_mono e "xinu_gui_add_slider"      (TFun ([TInt;TInt;TInt;TInt;TInt;TInt;TInt;TInt;TInt;TInt;TInt;TString], TUnit));
add_mono e "xinu_gui_slider_value"    (TFun ([TInt], TInt));
add_mono e "xinu_gui_set_fork_free"   (TFun ([TInt], TUnit));
add_mono e "xinu_gui_set_fork_held"   (TFun ([TInt; TInt], TUnit));
add_mono e "xinu_gui_set_phil"        (TFun ([TInt; TInt], TUnit));
add_mono e "xinu_gui_set_actor"       (TFun ([TInt; TInt; TInt], TUnit));
add_mono e "xinu_gui_buf_setup"       (TFun ([TInt; TInt; TInt], TUnit));
add_mono e "xinu_gui_buf_put"         (TFun ([TInt], TUnit));
add_mono e "xinu_gui_buf_take"        (TFun ([TInt], TUnit));
(* 関係演算子の int/float 混在 *)
let add_mxr1 f = add_mono e f (TFun ([TInt; TFloat], TBool)) in
List.iter add_mxr1 [ ">"; "<"; "<="; ">="; "=="; "!=" ];
let add_mxr2 f = add_mono e f (TFun ([TFloat; TInt], TBool)) in
List.iter add_mxr2 [ ">"; "<"; "<="; ">="; "=="; "!=" ];
\end{lstlisting}

\section{サンプルプログラムの修正}

ツールチェーン側を直した上で、3 サンプルにも\textbf{ 最小限の型注釈}を
入れた.いずれもセマンティクス変更は無く、AIPL の typed-field 構文
\texttt{var name: T = init;} を用いた\textbf{ ドキュメント的な追加}である.

\subsection{Rotate4LinesXinu.abcl}

\paragraph{変更点} (1) フィールド全てに型注釈を付与、
(2) \texttt{Spinner.init} の \texttt{sp: float, mx: float, my: float} を
\texttt{sp: int, mx: int, my: int} へ
(呼び出し側はもともと整数を渡していたため)、
(3) \texttt{Controller} のフィールド \texttt{s1..s4} を
空文字列ではなく \texttt{any} 注釈の整数 0 で初期化.

\begin{lstlisting}[language=abcl,caption={Rotate4LinesXinu.abcl 主な差分}]
class Spinner {
  var idx: int = 0;
  var angle: float = 0.;
  var speed: int = 5;
  var cx: int = 0;
  var cy: int = 0;
  var radius: int = 60;
  var running: int = 1;
  var rr: int = 200;
  var gg: int = 220;
  var bb: int = 255;

  method init(i: int, sp: int, mx: int, my: int,
              r: int, g: int, b: int) {
    ...
  }
}

class Controller {
  var s1: any = 0; var s2: any = 0;
  var s3: any = 0; var s4: any = 0;
  ...
}
\end{lstlisting}

\subsection{Philosophers5Xinu.abcl}

\paragraph{変更点}
\texttt{Fork.idx/held/holder} を \texttt{int} に明示、
\texttt{Philosopher} は actor 参照を保持する \texttt{left\_fork,
right\_fork} を \texttt{: any} 注釈、
\texttt{Controller} の \texttt{p0..p4, f0..f4} を \texttt{: any}.

\begin{lstlisting}[language=abcl,caption={Philosophers5Xinu.abcl actor 参照フィールド}]
class Philosopher {
  var idx: int       = 0;
  var state: int     = 0;
  var counter: int   = 0;
  var left_idx: int  = 0;
  var right_idx: int = 0;
  var left_fork:  any = 0;  // actor ref
  var right_fork: any = 0;  // actor ref
  var running: int   = 0;
  ...
}

class Controller {
  var p0: any = 0; var p1: any = 0; var p2: any = 0;
  var p3: any = 0; var p4: any = 0;
  var f0: any = 0; var f1: any = 0; var f2: any = 0;
  var f3: any = 0; var f4: any = 0;
  ...
}
\end{lstlisting}

\subsection{BoundedBufferXinu.abcl}

\paragraph{変更点}
\texttt{Buffer.capacity/count} を \texttt{int} 明示、
\texttt{Producer}/\texttt{Consumer} の \texttt{buffer} フィールドを
\texttt{: any} に
(actor 参照 \texttt{new Buffer(20)} を保持するため)、
\texttt{Controller} の \texttt{buf, p0..p2, c0..c2} を \texttt{: any}.

\begin{lstlisting}[language=abcl,caption={BoundedBufferXinu.abcl Producer のフィールド}]
class Producer {
  var idx: int     = 0;
  var state: int   = 0;
  var counter: int = 0;
  var buffer: any  = 0;   // Buffer actor ref
  var running: int = 0;
  ...
}

class Controller {
  var buf: any = 0;
  var p0: any = 0; var p1: any = 0; var p2: any = 0;
  var c0: any = 0; var c1: any = 0; var c2: any = 0;
  ...
}
\end{lstlisting}

\section{検証結果}

\subsection{R0 phase の acceptance — 5/5 PASS}

\begin{table}[h]
\centering
\small
\begin{tabular}{lccc}
\toprule
Assertion & R4L & P5 & BB \\
\midrule
(1) typecheck pass               & \checkmark & \checkmark & \checkmark \\
(2) C 生成成功 (\texttt{aipl2c}) & \checkmark & \checkmark & \checkmark \\
(3) Xinu kernel link 成功        & \checkmark & \checkmark & \checkmark \\
(4) QEMU \texttt{xsh\$} 到達     & \checkmark & \checkmark & \checkmark \\
(5) panic / FATAL 無し           & \checkmark & \checkmark & \checkmark \\
\bottomrule
\end{tabular}
\caption{R0 — 5 assertion × 3 sample = 15、全 PASS}
\end{table}

\subsection{後続 phase への波及}

R0 修正は\textbf{ AIPL 型システム全体に有効}な変更だったため、
後続 R1 / P1 phase でも追加の修復は不要だった:

\begin{table}[h]
\centering
\small
\begin{tabular}{lcc}
\toprule
Phase & Assertion 数 & 結果 \\
\midrule
R0 (型チェック修復)        & 5 & 5 PASS / 0 FAIL \\
R1 (smoke 自動化)           & 9 & 9 PASS / 0 FAIL \\
P1 (actor=Xinu thread)      & 12 & 12 PASS / 0 FAIL \\
\midrule
\textbf{合計}               & \textbf{26} & \textbf{26 PASS} \\
\bottomrule
\end{tabular}
\caption{Round 1 進捗 (2026-05-20)}
\end{table}

\subsection{Backward compat}

非 Xinu サンプル 8 件
(\texttt{Counter}, \texttt{PingPong}, \texttt{TypedDemo}, \texttt{Hello},
 \texttt{Generics}, \texttt{Philosophers5}, \texttt{Phase15\_Owned},
 \texttt{bounded\_buffer}) も全件 PASS で、特に\texttt{
TypedDemo.abcl} 自体が R0 修正によって\textbf{ 初めて型検査を
通るようになった}.

\subsection{コミット履歴}

\begin{lstlisting}[language=]
f5d902c  xinu-razpi-aipl P1: verify 1:1 actor->Xinu-thread mapping
a69308c  xinu-razpi-aipl R1: smoke automation with [aipl] markers
e80769e  xinu-razpi-aipl R0: restore typecheck for 3 Xinu samples   <-- 本報告
181ba26  xinu-razpi-aipl: add .ga.json (top-level) + .aipl orchestrator
dd35f63  xinu-razpi-aipl Round 1: design + MAP-Elites evolution
\end{lstlisting}

外部リポジトリ \texttt{xinu-raz/xinu} 側にも commit \texttt{5ee95b6}
(``Conditional AIPL auto-start via -DAIPL\_AUTOSTART'') を入れたが、
これは R1 smoke の都合上の auto-start hook であり、R0 とは独立.

\section{結論}

R0 phase は単なる「3 サンプルの型注釈を足す」作業に見えたが、
実際には AIPL 型システム本体に\textbf{ 構造的な不具合 3 つ}が
潜んでいた:

\begin{enumerate}[label=(\arabic*)]
  \item \texttt{infer.ml::check\_decl Class} がクラスフィールドの
        \texttt{TypedVarDecl} を無視していた
  \item \texttt{aipl2c.ml} のパイプラインが \texttt{normalize\_program} を
        呼んでおらず、 codegen 側の field 反復が \texttt{VarDecl} 限定
        だった
  \item \texttt{typing\_env.ml::prelude} に Xinu GUI 表面と
        int/float 混在の関係演算子が欠落していた
\end{enumerate}

これらを直したことで、3 Xinu サンプルだけでなく\textbf{ AIPL の
typed-field を持つ全プログラム}が再び型検査を通るようになり、
全 codegen バックエンド (Xinu / C / Pony / Erlang / Go / Prolog /
Python / OpenMP / LLVM) でも整合的な C を生成するようになった.
3 サンプルのソース側の修正は\textbf{ \emph{付随的}な型注釈追加}に
留まり、ロジックには一切手を加えていない.

\medskip

\noindent
これにより Round~1 の本丸 phase (P2-P4 / F1-F4 / G1-G5 / N1-N4) が
\textbf{ ツールチェーン側の足場が固まった状態で}着手できるように
なった.

\end{document}
